\subsection{滤子的比较}

定义 2. 给定同一集合 X 上的两个滤子 \(\mathfrak{F}\)、\(\mathfrak{F}'\)，称 \(\mathfrak{F}'\) 细于 \(\mathfrak{F}\)，或称 \(\mathfrak{F}\) 粗于 \(\mathfrak{F}'\)，如果 \(\mathfrak{F} \subset \mathfrak{F}'\)。若还有 \(\mathfrak{F} \neq \mathfrak{F}'\)，则称 \(\mathfrak{F}'\) 严格细于 \(\mathfrak{F}\)，或称 \(\mathfrak{F}\) 严格粗于 \(\mathfrak{F}'\)。

若两个滤子中一个细于另一个，则称这两个滤子可比较。X 上所有滤子所成之集按关系“F 粗于 F”排序；这个关系是由 \(\mathfrak{P}(\mathfrak{P}(\mathbf{X}))\) 中的包含关系诱导的。

设 \((\mathfrak{F}_t)_{\iota \in I}\) 是集合 \(X\) 上滤子的任一非空族（因而 X 必非空）；则集合

\[
\mathfrak {F} = \bigcap_ {i \in I} \mathfrak {F} _ {i}
\]

满足公理 \((\mathbf{F}_{\mathbf{I}})\)、\((\mathbf{F}_{\mathbf{II}})\) 与 \((\mathbf{F}_{\mathbf{III}})\)，从而是一个滤子；F 称为滤子族 \((\mathfrak{F}_{i})_{i\in\mathbf{I}}\) 的交，并且显然是 X 上所有滤子之有序集中诸 \(F_{i}\) 所成之集的下确界。

由单独一个集合 X 组成的滤子是 X 上所有滤子之有序集的最小元。我们将在第 4 目看到，若 X 的元素多于一个，则 X 上所有滤子所成之集没有最大元。

给定集合 X 的一个子集族 \(\mathfrak{G}\)，我们来考察 X 上是否存在包含 \(\mathfrak{G}\) 的滤子。若这样的滤子存在，则由 \((\mathrm{F}_{\mathrm{II}})\)，它也包含集 \(\mathfrak{G}'\)（\(\mathfrak{G}\) 中集合的有限交所成之集，包括 X，即 \(\mathfrak{G}\) 的空子集之交）；因此这样的滤子存在的一个必要条件是 X 的空子集不属于 \(\mathfrak{G}'\)。这个条件也是充分的，因为由 \((\mathrm{F}_{\mathrm{I}})\)，任何包含 \(\mathfrak{G}'\) 的滤子也包含集 \(\mathfrak{G}''\)（X 中包含 \(\mathfrak{G}'\) 中某个集合的子集所成之集）。而 \(\mathfrak{G}''\) 显然满足 \((\mathrm{F}_{\mathrm{I}})\)；它满足 \((\mathrm{F}_{\mathrm{II}})\)（由 \(\mathfrak{G}'\) 的定义）；最后，它满足 \((\mathrm{F}_{\mathrm{III}})\)，因为 X 的空子集不属于 \(\mathfrak{G}'\)。于是 \(\mathfrak{G}''\) 是包含 \(\mathfrak{G}\) 的最粗滤子，而且我们已证明：

命题 1. X 上存在包含 X 的子集所成之集 \(\mathfrak{G}\) 的滤子的充分必要条件是：\(\mathfrak{G}\) 的任何有限子集的交都非空。

上面定义的滤子 \(\mathfrak{G}''\) 称为由 \(\mathfrak{G}\) 生成的滤子，而 \(\mathfrak{G}\) 称为 \(\mathfrak{G}''\) 的一个子基。

例. 设 \(\mathfrak{S}\) 是集合 X 的子集所成的任一集，并设 \(\mathfrak{C}\) 是由 \(\mathfrak{S}\) 生成的 X 上的拓扑（\(\S\) 2，第 3 目，例 II）。由于 \(\mathfrak{S}\) 中集合的有限交所成之集是 \(\mathfrak{C}\) 的一个基，故由上面命题 1 的证明以及 \(\S\) 1 第 3 目命题 3 可知，对每个 \(x \in X\)，\(x\) 关于 \(\mathcal{C}\) 的邻域滤子由集合 \(\mathfrak{S}(x)\)（\(\mathfrak{S}\) 中包含 \(x\) 的集合所成之集）生成。

推论 1. 设 \(\mathfrak{F}\) 是集合 X 上的一个滤子，A 是 X 的一个子集。则存在滤子 \(\mathfrak{F}'\)，它细于 \(\mathfrak{F}\) 且使得 \(A \in \mathfrak{F}'\)，当且仅当 A 与 \(\mathfrak{F}\) 的所有集合相交。

推论 2. 滤子所成之集 \(\Phi\)（非空集合 \(\mathbf{X}\) 上的）在 \(\mathbf{X}\) 上所有滤子所成之集中有上确界，当且仅当对所有有限序列 \((\mathfrak{F}_i)_{1\leqslant i\leqslant n}\)（\(\Phi\) 的元素的）以及所有 \(\mathrm{A}_i\in \mathfrak{F}_i\)（\(1\leqslant i\leqslant n\)），交 \(\mathrm{A}_1\cap \dots \cap \mathrm{A}_n\) 都非空。

因为这一条件表明并 \(\mathfrak{G}\)——诸滤子 \(\mathfrak{F} \in \Phi\) 之并——满足命题 1 的条件。

推论 3. 非空集合 X 上所有滤子所成之有序集是归纳的。

因为滤子的每个线性有序集 \(\Phi\)（\(X\) 上的）都满足命题 1 的推论 2 的条件，这是由于按假设诸集合 \(A_{i}\) 都属于同一个 \(\mathfrak{F}_{j}\)，于是可以应用 \((\mathbf{F}_{\mathrm{II}})\)。

